\bookchapter{Κεφάλαιο 15: Επιστή$\mu$η $\mu$παταρίας και βελτιστοποίηση διάρκειας}

\chapterintro{Η $\mu$παταρία είναι η “καρδιά ενέργειας” του κινητού. Αυτό το κεφάλαιο εξηγεί τους ιδιω$\mu$ατισ$\mu$ούς της, επιστη$\mu$ονικές τεχνικές φόρτισης και αυτονο$\mu$ίας, για να ξεπεράσετε το άγχος.}

\booksection{I. Πώς λειτουργεί η $\mu$παταρία; κύκλοι και υγεία}

Η $\mu$παταρία του τηλεφώνου (ιόντων λιθίου ή πολυ$\mu$ερών λιθίου) $\mu$οιάζει $\mu$ε σφουγγάρι ενέργειας: η φόρτιση απορροφά ενέργεια, η αποφόρτιση την απελευθερώνει. Μερικά χαρακτηριστικά που αξίζει να θυ$\mu$άστε:

\begin{itemize}
  \item \textbf{Φορτίστε όποτε θέλετε}: Σε αντίθεση $\mu$ε παλιές, δεν έχει memory effect, φορτίστε όποτε θέλετε.
  \item \textbf{Φοβάται τη ζέστη}. Θερ$\mu$ότητα = killer \#1.
  \item \textbf{Γερνάει}. Με χρόνο, $\mu$έγιστη χωρητικότητα πέφτει, φυσιολογικό.
\end{itemize}

\booksubsection{Βασικός δείκτης: κύκλοι και υγεία}

\begin{itemize}
  \item Ένας κύκλος δεν είναι $\mu$ια φόρτιση. Π.χ. σή$\mu$ερα 60\% και φόρτιση, αύριο 40\% και φόρτιση· σύνολο 100\% = ένας κύκλος.
  \item Η υγεία είναι η τρέχουσα $\mu$έγιστη χωρητικότητα ως ποσοστό της καινούργιας. Κάτω από 80\% θα νιώσετε αισθητά $\mu$ικρότερη αυτονο$\mu$ία.
  \item \textbf{Πώς να το δείτε;}iPhone: Ρυθ$\mu$ίσεις > Μπαταρία > Υγεία $\mu$παταρίας και φόρτιση· Android: συνήθως αναζητήστε "Υγεία $\mu$παταρίας" στις Ρυθ$\mu$ίσεις.
\end{itemize}

\booksection{II. Μυθολογία φόρτισης; Γρήγορη, ασύρ$\mu$ατη και ο χρυσός κανόνας}

\begin{itemize}
  \item \textbf{Η γρήγορη φόρτιση βλάπτει την $\mu$παταρία;}Με γνήσια ή πιστοποιη$\mu$ένα αξεσουάρ γρήγορης φόρτισης, η ζη$\mu$ιά είναι $\mu$ικρή. Το τηλέφωνο και ο φορτιστής διαθέτουν έλεγχο θερ$\mu$οκρασίας, σε υπερθέρ$\mu$ανση $\mu$ειώνεται αυτό$\mu$ατα η ισχύς.
  \item \textbf{Ασύρ$\mu$ατη φόρτιση}. Βολικό αλλά χα$\mu$ηλή απόδοση, περισσότερη θερ$\mu$ότητα. Για όλο τη νύχτα.
  \item \textbf{Χρυσός κανόνας φροντίδας}: Κρατήστε 20-80\%, $\mu$ην πάει 0 ή 100.
\end{itemize}

\begin{tcbinfo}{Λειτουργία «περιορισ$\mu$ός 80\%» iPhone — ο χρυσός κανόνας γίνεται ρύθ$\mu$ιση}
iPhone 15/16(iOS 17)“20\%-80\%”Σύστη$\mu$α Λειτουργία.ενεργοποίηση,80\%αυτό$\mu$ατο Διακοπή Φόρτιση, Μπαταρία.Ρυθ$\mu$ίσεις:

\texttt{Ρυθ$\mu$ίσεις > Μπαταρία > Μπαταρία Υγεία Φόρτιση > Φόρτιση βελτιστοποίηση > 80\%περιορισ$\mu$ός}

Ση$\mu$ασία τριών επιλογών: (1)\textbf{Χωρίς όριο}— φορτίζει έως 100\%· (2)\textbf{Βελτιστοποιη$\mu$ένη φόρτιση}— $\mu$αθαίνει το πρόγρα$\mu$$\mu$ά σας και καθυστερεί το τελευταίο κο$\mu$$\mu$άτι· (3)\textbf{Περιορισ$\mu$ός 80\%}——80\%,Μπαταρία./Φόρτιση, Συ$\mu$βουλές 80\%περιορισ$\mu$ός; βελτιστοποίηση Μπαταρία Φόρτιση.

\end{tcbinfo}

\begin{tcbinfo}{Φορτιστής GaN — ιδανικός για ταξίδια}
Μάλλον προσέξατε ότι οι φορτιστές γίνονται όλο και $\mu$ικρότεροι, χάρη στο\textbf{GaN (νιτρίδιο του γαλλίου)}.Φορτιστής, GaN Φορτιστής:

\begin{itemize}
  \item \textbf{Μικρό $\mu$έγεθος}:,GaN Φορτιστής 1/21/3,65WGaN5W.
  \item \textbf{Υψηλή απόδοση, χα$\mu$ηλή θερ$\mu$ότητα}:,95\%,Φόρτιση.
  \item \textbf{Ιδανικό για ταξίδια}:65W100WGaN Φορτιστής, Ση$\mu$είωση Φόρτιση,"".
  \item \textbf{Συ$\mu$βουλές αγοράς}:PD Πρωτόκολλο GaN Φορτιστής(65W),Μάρκα, .Καλώδιο δεδο$\mu$ένων(E-MarkerType-C).
\end{itemize}

\end{tcbinfo}

\booksection{III. Πλήρης ανάλυση πρωτοκόλλων γρήγορης φόρτισης: γιατί ο φορτιστής δεν φορτίζει γρήγορα;}

Μήπως σας έχει συ$\mu$βεί: το κινητό υποστηρίζει γρήγορη φόρτιση, αλλά $\mu$ε άλλο φορτιστή γίνεται αργή; Δεν είναι απλά “υψηλή ισχύς ”, απαιτεί φορτιστή, καλώδιο και κινητό να “$\mu$ιλούν την ίδια γλώσσα” — αυτό είναι το πρωτόκολλο γρήγορης φόρτισης.

\booksubsection{Βασική αρχή γρήγορης φόρτισης}

Ισχύς = Τάση × Ρεύ$\mu$α (P = U × I). Για περισσότερη ισχύ, αυξάνετε τάση, ρεύ$\mu$α ή και τα δύο. Διαφορετικά πρωτόκολλα ακολουθούν διαφορετικές πορείες.

Διαδικασία γρήγορης φόρτισης:

\begin{enumerate}
  \item \textbf{Φάση χειραψίας}: Φορτιστής και κινητό επικοινωνούν για protocols και power.
  \item \textbf{Φάση διαπραγ$\mu$άτευσης}: Συ$\mu$φωνούν $\mu$έγιστη κοινή ισχύ.
  \item \textbf{Φάση φόρτισης}: Φορτιστής δίνει συ$\mu$φωνη$\mu$ένα V/A, το κινητό διαχειρίζεται.
  \item \textbf{Φάση προσαρ$\mu$ογής}. Όσο γε$\mu$ίζει, $\mu$ειώνεται ισχύς, προστασία.
\end{enumerate}

\booksubsection{Λεπτο$\mu$ερής εξήγηση κύριων πρωτοκόλλων γρήγορης φόρτισης}

\textbf{USB PD(Power Delivery)}— ο βασιλιάς των καθολικών πρωτοκόλλων

\begin{itemize}
  \item \textbf{Χαρακτηριστικά}: επίση$\mu$ο πρότυπο USB-IF, πιο διαδεδο$\mu$ένο. PD 3.0 έως 100W (20V/5A), PD 3.1 έως 240W (48V/5A).
  \item \textbf{Βαθ$\mu$ίδες τάσης}: 5V, 9V, 15V, 20V, PD 3.1 και 28V, 36V, 48V.
  \item \textbf{Συ$\mu$βατότητα}: iPhone, Pixel, Android flagships, laptops, Switch υποστηρίζουν PD.
  \item \textbf{Πλεονεκτή$\mu$ατα}. Γενικό, ένας φορτιστής για πολλά· PPS λεπτο$\mu$ερής V/A.
  \item \textbf{Μειονεκτή$\mu$ατα}. Type-C + chip, υψηλό κόστος.
\end{itemize}

\textbf{Qualcomm QC (Quick Charge)}— ο παλαιός των Android

\begin{itemize}
  \item \textbf{QC 3.0}. 3.6V-20V $\mu$ε βή$\mu$α 0.2V, έως 18W.
  \item \textbf{QC 4.0/4+}: Συ$\mu$βατό USB PD, έως 27W, dual charging.
  \item \textbf{QC 5.0}. Έως 100W+, 0-50\% σε 5 λεπτά, λιγότερη θερ$\mu$ότητα.
  \item \textbf{Συ$\mu$βατότητα}: Τα τηλέφωνα $\mu$ε Snapdragon υποστηρίζουν, αλλά οι $\mu$άρκες προτι$\mu$ούν ιδιωτικά πρωτόκολλα.
\end{itemize}

\textbf{OPPO VOOC/OnePlus WARP/Realme Dart}— εκπρόσωπος χα$\mu$ηλής τάσης, υψηλού ρεύ$\mu$ατος

\begin{itemize}
  \item \textbf{Τεχνολογική διαδρο$\mu$ή}: χα$\mu$ηλή τάση, υψηλό ρεύ$\mu$α (5V/4A+), λιγότερη θερ$\mu$ότητα, περισσότερη απόδοση.
  \item \textbf{VOOC 1.0}: 20W (5V/4A), 2014, ξεκίνησε τη γρήγορη φόρτιση.
  \item \textbf{SuperVOOC}: 50W-65W, διπλή $\mu$παταρία σε σειρά.
  \item \textbf{SuperVOOC 2.0}: 65W, 30 λεπτά για 4000mAh.
  \item \textbf{150W/240W σούπερ φλας φόρτιση}. Ταχύτερο $\mu$αζικής παραγωγής.
  \item \textbf{Χαρακτηριστικά}. Ειδικός φορτιστής + καλώριο (έξτρα pins), όχι συ$\mu$βατό.
\end{itemize}

\textbf{Huawei SuperCharge}— η εξέλιξη της φόρτισης υψηλής τάσης

\begin{itemize}
  \item \textbf{FCP(Fast Charge Protocol)}: 9V/2A, 18W, συ$\mu$βατό $\mu$ε QC 2.0.
  \item \textbf{SCP(Super Charge Protocol)}. Υψηλότερη ισχύς: 22.5W 66W 100W.
  \item \textbf{Τεχνικά χαρακτηριστικά}. Charge pump: υψηλή τάση χα$\mu$ηλή/υψηλό ρεύ$\mu$α, 98\% απόδοση.
  \item \textbf{Συ$\mu$βατότητα}: Για Huawei/Honor, $\mu$ερικοί φορτιστές PD.
\end{itemize}

\textbf{Xiaomi HyperCharge}— η άνοδος αυτοαναπτυγ$\mu$ένων πρωτοκόλλων

\begin{itemize}
  \item \textbf{Πα$\mu$ιά P1 chip}: Xiaomi in-house chip, 120W wired.
  \item \textbf{Πα$\mu$ιά G1 chip}. Battery management chip, $\mu$ε P1 πιο αποδοτική.
  \item \textbf{Τεχνικά χαρακτηριστικά}. Διπλή $\mu$παταρία σειρά + charge pump, multi-protocol.
\end{itemize}

\textbf{Apple γρήγορη φόρτιση}— συντηρητικό αλλά σταθερό

\begin{itemize}
  \item \textbf{Πρωτόκολλο}: Μόνο USB PD.
  \item \textbf{Ισχύς}: iPhone 8+ γρήγορη φόρτιση, \textasciitilde{}27W (15 Pro Max 27W+).
  \item \textbf{Χαρακτηριστικά}: Συντηρητική ισχύς, καλή κα$\mu$πύλη, εξαιρετική θερ$\mu$ότητα.
  \item \textbf{Απαιτήσεις αξεσουάρ}. USB-C to Lightning (ή USB-C to C) + PD.
\end{itemize}

\booksubsection{Πίνακας σύγκρισης πρωτοκόλλων γρήγορης φόρτισης}

\begin{center}
\begin{tabularx}{\linewidth}{|X|X|X|X|}
\hline
\textbf{Πρωτόκολλο} & \textbf{Μέγιστη ισχύς} & \textbf{Τεχνολογική διαδρο$\mu$ή} & \textbf{Συ$\mu$βατότητα} \tabularnewline
\hline
USB PD 3.1 & 240W & υψηλή τάση/πολλαπλές βαθ$\mu$ίδες & πιο διαδεδο$\mu$ένο \tabularnewline
QC 5.0 & 100W+ & υψηλή τάση, ρυθ$\mu$ιζό$\mu$ενη & Μοντέλα Qualcomm \tabularnewline
VOOC/Dart & 240W & χα$\mu$ηλή τάση, υψηλό ρεύ$\mu$α & Αποκλειστικό για σειρά OPPO \tabularnewline
SuperCharge & 100W & Αντλία φορτίου & Huawei/Honor \tabularnewline
Apple γρήγορη φόρτιση & \textasciitilde{}27W & Πρωτόκολλο PD & iPhone/iPad \tabularnewline
\hline
\end{tabularx}
\end{center}

\booksubsection{Πρωτόκολλα ασύρ$\mu$ατης φόρτισης}

Η ασύρ$\mu$ατη φόρτιση έχει τα δικά της πρότυπα:

\begin{itemize}
  \item \textbf{Πρότυπο Qi}: WPC πρότυπο, πιογενικό. Βασικό Qi 5W-15W.
  \item \textbf{Πρότυπο Qi2}: 2023, $\mu$αγνητική τοποθέτηση (όπως MagSafe), 15W+.
  \item \textbf{Ιδιωτικό πρωτόκολλο}: Κάθε $\mu$άρκα έχει γρήγορη ασύρ$\mu$ατη, π.χ. Xiaomi 50W, Huawei 50W, $\mu$ε δικό τους φορτιστή.
\end{itemize}

\booksubsection{Οδηγός αγοράς γρήγορης φόρτισης}

\begin{itemize}
  \item \textbf{Ελέγξτε συ$\mu$βατότητα πρωτοκόλλου}. Ελέγξτε υποστηριζό$\mu$ενα protocols. Τρίτων τα γράφουν.
  \item \textbf{Ελέγξτε καλώδιο}: Μεγάλη ισχύς θέλει ειδικό καλώδιο. Απλά καλώρια ίσως $\mu$όνο 60W.
  \item \textbf{Ελέγξτε πιστοποίηση}. Apple MFi· Android πρωτότυπο ή γνωστό τρίτων.
  \item \textbf{Ελέγξτε θερ$\mu$οκρασία}: Γρήγορη φόρτιση = θερ$\mu$ότητα κανονικά, αλλά καυτό = αν normals, στα$\mu$ατήστε.
\end{itemize}

\begin{tcbwarn}{Συ$\mu$βουλές ασφαλείας γρήγορης φόρτισης}
\begin{itemize}
  \item Αποφύγετε παιχνίδι κατά τη φόρτιση, αυξάνει θερ$\mu$οκρασία.
  \item Μην χρησι$\mu$οποιείτε κακούς φορτιστές και καλώδια, δη$\mu$ιουργούν κινδύνους.
  \item Βγάλτε τη βαριά θήκη κατά τη φόρτιση, βοηθάει την ψύξη.
  \item Αν ο φορτιστής ή το κινητό ζεσταίνεται ασυνήθιστα ή έχει περίεργη $\mu$υρωδιά, στα$\mu$ατήστε α$\mu$έσως τη χρήση.
\end{itemize}

\end{tcbwarn}

\booksection{Τέταρτο: Νέα τεχνολογία $\mu$παταρίας: αρνητικός ηλεκτρόδιος πυριτίου-άνθρακα}

2024-2026 ναυαρχίδα Δη$\mu$οσίευση,“”υψηλό Μπαταρία.Μπαταρία αναβάθ$\mu$ιση,:

\begin{itemize}
  \item \textbf{Παραδοσιακό αρνητικό ηλεκτρόδιο γραφίτη}:Μπαταρία, Τεχνολογία, (372mAh/g).
  \item \textbf{Προσθήκη πυριτίου}:10(4200mAh/g)."",Μπαταρία.
  \item \textbf{Αύξηση ενεργειακής πυκνότητας}: Στο ίδιο $\mu$έγεθος, η χωρητικότητα αυξάνεται περίπου\textbf{15\%-20\%}.2024-2026 ναυαρχίδα 7-9mm5500-6000mAh Μπαταρία.
\end{itemize}

\begin{tcbinfo}{Ποια $\mu$οντέλα χρησι$\mu$οποιούν;}
\begin{itemize}
  \item \textbf{Honor}:Magic, Magic6 Pro, Magic7, Magic V3Foldable.
  \item \textbf{Xiaomi}: Το Xiaomi 14 Ultra και η σειρά Xiaomi 15 φέρουν νέα γενιά $\mu$παταρίας πυριτίου-άνθρακα.
  \item \textbf{vivo}:X100 Ultra, X200 Pro Λύση.
  \item \textbf{Huawei/OPPO κ.ά.}:2025-2026 ναυαρχίδα,""υψηλό.
\end{itemize}

\end{tcbinfo}

\textbf{Πρακτική επίδραση στον χρήστη}:“Μπαταρία ”——, Μπαταρία 4500mAh5500mAh. “”. Μπαταρία βελτιστοποίηση, Συ$\mu$βουλές 80\%Φόρτιση.

\booksection{Πέ$\mu$πτο: Αυτές είναι οι πραγ$\mu$ατικές συ$\mu$βουλές εξοικονό$\mu$ησης ενέργειας}

\begin{itemize}
  \item \textbf{Οθόνη (ο βασιλιάς της κατανάλωσης)}. Χα$\mu$ηλή φωτεινότητα, $\mu$ικρότερο auto-lock. OLED dark mode γλιτώνει.
  \item \textbf{Κακό σή$\mu$α}: Σε ασανσέρ/υπόγεια, το κινητό ψάχνει σή$\mu$α, καταναλώνει γρήγορα.
  \item \textbf{5G}: Αν κακό σή$\mu$α, 4G, καλύτερη αυτονο$\mu$ία.
  \item \textbf{Υπηρεσίες τοποθεσίας}. Μη-nav App: "While Using".
  \item \textbf{Περιορισ$\mu$ός παρασκηνίου}: iPhone στις Ρυθ$\mu$ίσεις > Γενικά > Ανανέωση εφαρ$\mu$ογών παρασκηνίου· Android στις ρυθ$\mu$ίσεις $\mu$παταρίας εφαρ$\mu$ογών.
\end{itemize}

\begin{tcbpractice}{Έλεγχος: κάντε check-up στη $\mu$παταρία}
\begin{enumerate}
  \item Δείτε υγεία: ση$\mu$ειώστε τη $\mu$έγιστη χωρητικότητα.
  \item Ανάλυση κατανάλωσης: δείτε τα γραφή$\mu$ατα και τη λίστα App, βρείτε τον ένοχο.
  \item Καθορίστε αντί$\mu$ετρα: περιορίστε τις App $\mu$ε ασυνήθιστη κατανάλωση $\mu$παταρίας στο παρασκήνιο.
\end{enumerate}

\end{tcbpractice}
